\section{Prepare a cumulative observation}
Use \texttt{drafts/matmul-construction/main.asm} as the fresh starter. The scalar
entry is selected initially. Each marked region contains a temporary branch and a
\texttt{nop} that return through the supplied restoring exit. Replace both lines
with the corresponding increment; keep surrounding labels and markers. Later code
is present as safe regions, so each cumulative program assembles before it is a
complete multiplication function.

The starter uses \texttt{.word 0x5a5aa5a5:544} to initialize 544 words with the same sentinel pattern. The repetition count describes the backing arena, not the active matrix size.

The input recipes accompany every edit. Use \texttt{prepare.py} with a recipe ID,
your current cumulative source and a separate output path. The preparation changes
only the input words, shape and caller method. It refuses to overwrite an existing
observation file or the source being prepared. Start from a program that retains
the default shape and method lines; the recipe selects the function for that run.

The caller initializes an aligned stack and loads the three matrix pointers and
the descriptor pointer. It stops at \texttt{matmul\_done} immediately after the
selected function returns. The result count is still in \texttt{v0}; the caller's
\texttt{completed} word is still \(-1\). Continue to \texttt{results\_ready} to
see that word updated. The function's supplied exit has already restored saved
registers. Inspect the documented temporary registers for partial construction
state, and inspect memory for any stored partial sums.

The verification runs each observation through this caller and through a direct
function call with independent nonzero offsets in the input and output arenas.
It checks the full pointer values, exact stored bytes, preservation requirements,
logical access intervals and initialized partial-sum reads. An unfinished program
can return the correct cell count while its product remains incomplete. Final
behavioral checks are a separate requirement from these construction observations.

The code uses \texttt{s0}, \texttt{s1}, \texttt{s2}, \texttt{s3} for \(M,K,N,T\),
\texttt{s4}, \texttt{s5}, \texttt{s6} for the full matrix bases, \texttt{s7} for
row-tile start \(I\), and \texttt{fp} for column-tile start \(J\).
The tiled reduction start \(L\) is in \texttt{t0}; current indices \(i,j,k\)
are in \texttt{t1}, \texttt{t2}, \texttt{t3}. This is the illustrated allocation.
An equivalent free-writing solution may use different internal state.
